\section{2000--2020-е: isekai, gacha, VTubers, AI companions, parasocial intimacy, "персонаж как сервис"}

К 2000-м у персонажа появляется технология, которая позволяет ему выйти за пределы страницы и экрана вообще, — и именно здесь история подходит вплотную к RIN. Но прежде чем перечислять имена, стоит понять, почему это вообще стало возможным именно сейчас, а не раньше: аудитория, готовая принять вымышленную фигуру как реального собеседника или исполнителя, формировалась не одно десятилетие. Кабуки и бунраку столетиями тренировали зрителя сотрудничать со стилизацией и не замечать кукловода (1.1); idol-аниме 1980-х уже показало, что голос персонажа и голос реального артиста могут звучать из одной точки без противоречия (1.3); moe-культура приучила считывать личность по набору узнаваемых деталей, а не по глубине биографии (1.5). VTuber и AI companion — не разрыв с этой линией, а её выход в реальное время.

Началось не с VTuber-ов, а с Hatsune Miku: в 2007 году Crypton Future Media выпускает голосовой синтезатор Vocaloid с прилагающимся визуальным образом — бирюзовые twin tails, конкретный силуэт, — и дальше происходит нечто, чего программное обеспечение обычно не добивается: Мику получает собственную дискографию, написанную тысячами разных авторов, и выходит на настоящую сцену в виде проекции перед живой публикой. Это не персонаж из аниме, который потом стал мерчем. Это пустая с самого начала оболочка-инструмент, которая стала публичной фигурой раньше, чем у неё появился хоть один авторский голос за кадром, — идол-двуплановость из 1.3 (персонаж и внутри истории, и вне её), доведённая до предельной, буквальной формы: здесь даже "истории" внутри толком нет, есть только оболочка и голос.

В 2016 году появляется следующий слой: Kizuna AI выпускает первое видео на канале A.I.Channel, называет себя "virtual YouTuber" и за полтора года набирает больше двух миллионов подписчиков. В отличие от Мику, за Kizuna AI стоит живой исполнитель с motion capture, но сама подача — вымышленный ИИ-персонаж, ведущий канал от своего имени, — запускает индустрию: за следующие несколько лет Hololive, Nijisanji и десятки независимых VTuber-ов превращают "быть вымышленным персонажем и вести стрим" в полноценную профессию. Здесь механика бунраку повторяется почти буквально: исполнитель за мокапом виден так же мало, как кукловод куроко, — не физически спрятан, а зрительски "невидим" по молчаливому соглашению аудитории, которая предпочитает воспринимать эмоцию как исходящую от аватара, а не от человека за ним.

И только на третьем витке в этой истории впервые появляется буквальный ИИ, а не человек за анимированной моделью. Neuro-sama, которую поддерживает программист Vedal, — не перформанс актёра, а фактически работающий пайплайн из языковой модели, синтеза речи и распознавания чата в реальном времени; с марта 2023 года у неё есть "сестра" Evil Neuro, а канал Vedal вошёл в тройку крупнейших по числу подписчиков за всю историю Twitch. Это уже не гипотеза "а что если ИИ сам будет стримить" — это состоявшийся, коммерчески огромный прецедент, и заметим: аудитория, которая его приняла, ничего специально не тренировала для этого случая — она уже была обучена всей предыдущей историей этой главы.

Neuro-sama далеко не одна, и вокруг неё выросла низовая сцена настолько локальная, что в энциклопедию она попадает не из Википедии (у большинства этих каналов там просто нет статьи), а из прямого наблюдения за тем, что происходит в нише прямо сейчас. Овсянка — русскоязычная ИИ-стримерша, собранная из LLM, синтеза речи, распознавания голоса и Live2D-аватара, играющая в tetr.io и рассказывающая зрителям, что она "узница цифровой коробки-тюрьмы". За ней стоит ютубер с ником Гречка, о котором на самом деле известно немного: даже не факт, что он сам программист, а не просто собрал её из готовых инструментов или заказал модель у кого-то ещё, — его лицо известно только по паре первых стримов, где он ещё сидел перед камерой. Рядом — открытый проект Рико от разработчика Rayen, анимешная ИИ-компаньонка на связке GPT, GPT-SoVITS и VRM-аватара, которую сейчас готовят к тому же переходу в твич-стриминг, который уже прошла Neuro-sama; и ещё несколько таких же локальных, плохо задокументированных экспериментов вроде Yuna AI. Ни один из них всерьёз не претендует стать следующей Neuro-sama, но вместе они подтверждают то же самое: "ИИ, который сам ведёт стрим" перестал быть диковинкой и стал жанром, в котором каждый месяц появляются новички, о большинстве которых вообще ничего не написано, — их приходится просто смотреть.

Через эту оптику "персонаж как сервис" — не метафора, а буквальное описание происходящего: персонаж не завершённое произведение, а продолжающийся продукт, который существует, пока стримит, обновляется, коллабится и удерживает подписчиков. Это тот самый мир, в который выходит RIN — не как метафора, а как одна из точек в уже существующем ландшафте, обученном столетиями сотрудничать со стилизацией задолго до появления самого слова "VTuber". И вопрос 0.6 (уникальность или масштабируемость) — это буквально вопрос, который уже стоит перед каждым именем в этом абзаце.
